\chapter{Deployment}
\label{ch:deployment}

\section{Deployment inventory}

The project deploys three distinct things through three different channels.

\begin{longtable}{p{3.6cm}p{3.2cm}p{4.2cm}p{3.2cm}}
\caption{Deployment inventory}\\
\toprule
\textbf{Artefact} & \textbf{Channel} & \textbf{Automation} & \textbf{Status} \\
\midrule

\texttt{@expo-snips/core} & npm public registry & \texttt{publish.yml} on \texttt{v*}
tags & IMPLEMENTED \\
\texttt{@expo-snips/cli} & npm public registry & \texttt{publish.yml} on \texttt{v*}
tags & IMPLEMENTED \\
Screen registry (JSON) & GitHub Pages & \texttt{deploy-registry.yml} on
\texttt{main}/\texttt{master} & IMPLEMENTED \\
VS Code extension (\texttt{.vsix}) & Manual / Marketplace & None & NOT AUTOMATED \\
\texttt{packages/registry} (npm) & --- & \texttt{private: true} & NOT PUBLISHED \\
\bottomrule
\end{longtable}

\section{Registry deployment}

The registry is a folder of static JSON: \texttt{index.json} plus one document
per recipe under \texttt{recipes/}. Its deployment is:

\begin{enumerate}
  \item Checkout the repository.
  \item \texttt{node scripts/build-registry.mjs} --- this \emph{deletes} the
        existing \texttt{registry/} directory, rebuilds it from
        \texttt{packages/registry/data} with every template inlined, and writes
        \texttt{index.json} plus \texttt{recipes/<id>.json}.
  \item Copy the built \texttt{registry/} folder into \texttt{public/registry/}.
  \item Publish \texttt{public/} to the \texttt{gh-pages} branch with
        \texttt{peaceiris/actions-gh-pages@v3}, using the workflow's
        \texttt{GITHUB\_TOKEN}.
\end{enumerate}

The workflow runs on every push to \texttt{main} or \texttt{master} and can be
triggered manually via \texttt{workflow\_dispatch}. It declares
\texttt{permissions: contents: write}, which the Pages action needs to push the
\texttt{gh-pages} branch.

\textbf{Client default.} Clients resolve a registry base from configuration; when
none is supplied, the CLI's built-in constant points at
\texttt{https://registry.expo-snips.dev} and the VS Code extension defaults to its
\emph{embedded} catalog (empty registry URL). The repository remote is
\texttt{github.com/unkown812/expo-snips}, and the deployed Pages site would be
served under that account's GitHub Pages domain.

\textbf{INFERRED:} the hard-coded default host
\texttt{registry.expo-snips.dev} is not configured anywhere in this repository
(no CNAME file, no DNS configuration). Whether it resolves to the GitHub Pages
deployment is \textbf{UNKNOWN} from the repository alone.

\section{npm package deployment}

\texttt{publish.yml} triggers on tags matching \texttt{v*} and publishes the two
public packages:

\begin{lstlisting}[language=YAML]
- run: pnpm install --frozen-lockfile
- run: pnpm build
- run: pnpm --filter @expo-snips/core publish --no-git-checks --access public
  env: { NODE_AUTH_TOKEN: ${{ secrets.NPM_TOKEN }} }
- run: pnpm --filter @expo-snips/cli publish --no-git-checks --access public
  env: { NODE_AUTH_TOKEN: ${{ secrets.NPM_TOKEN }} }
\end{lstlisting}

Both packages declare \texttt{"publishConfig": \{"access": "public"\}} and a
\texttt{files: ["dist"]} allow-list, so only compiled output is uploaded. The
CLI's build step first runs \texttt{scripts/prepare-registry.mjs}, which inlines
the catalog into \texttt{src/generated/registry-data.ts} before
\texttt{tsc} compiles it --- so the published npm package is self-contained and
works offline.

\textbf{Note:} the workflow publishes only \texttt{core} and \texttt{cli}. The
\texttt{registry} package is marked \texttt{private} and is never published.

\section{VS Code extension deployment}

\textbf{Status:} NOT AUTOMATED.

Packaging is manual:

\begin{lstlisting}[language=bash]
pnpm --filter expo-snips package   # esbuild + vsce package
\end{lstlisting}

This produces \texttt{expo-snippets-<version>.vsix}. Four packaged extensions
(0.1.0 through 0.1.3) are present in the working tree under
\texttt{packages/vscode-extension/} (they are excluded from git by a
\texttt{*.vsix} rule in \texttt{.gitignore}), which indicates the current
distribution practice: build locally and share or install the artefact manually
(``Install from VSIX''). There is no publish step to the VS Code Marketplace and
no workflow that builds or attaches the VSIX to a GitHub release.

\textbf{RECOMMENDED:} add a workflow that packages the extension on a tag and
either publishes to the Marketplace (\texttt{vsce publish}) or attaches the VSIX
to a GitHub release, so the artefact is reproducible rather than committed.

\section{Runtime requirements and prerequisites}

\begin{longtable}{p{4cm}p{10.2cm}}
\caption{Runtime prerequisites}\\
\toprule
\textbf{Component} & \textbf{Requirement} \\
\midrule

Node.js (CLI and build) & $\geq$ 18 (declared in \texttt{engines} and in the
esbuild \texttt{target: node18}). CI uses Node 20. \\
pnpm & 9 (workflow pins \texttt{version: 9}); 9+ recommended. \\
VS Code (extension) & $\geq$ 1.80.0 (declared in \texttt{engines.vscode}). \\
Target project & Any Expo / React Native project; detection reads
\texttt{package.json} and \texttt{tsconfig.json}. \\
Package manager (for deps) & Any of npm, pnpm, yarn or bun; the extension detects
which and the CLI delegates to the project's installer. \\
\bottomrule
\end{longtable}

\section{What is absent}

\begin{itemize}
  \item \textbf{No containerisation.} There is no \texttt{Dockerfile},
        \texttt{docker-compose} file or container registry configuration.
  \item \textbf{No infrastructure-as-code.} Hosting is GitHub Pages; there is no
        Terraform, Pulumi or cloud configuration.
  \item \textbf{No environment-specific build.} Builds do not branch on an
        environment; there is a single production shape.
  \item \textbf{No rollback procedure.} Nothing documents how to revert a bad
        registry publish. Because registry state is rebuilt from the repository,
        reverting the commit and re-running the workflow restores the previous
        state --- but this is INFERRED, not documented.
  \item \textbf{No release notes automation.} The extension has a manual
        CHANGELOG; the npm packages have no CHANGELOG.
\end{itemize}

\section{Deployment summary}

\begin{table}[H]
\centering
\caption{Deployment posture}
\begin{tabular}{p{4cm}p{10.2cm}}
\toprule
\textbf{Dimension} & \textbf{Assessment} \\
\midrule
Registry delivery & IMPLEMENTED and automated (GitHub Pages, static JSON). \\
npm delivery & IMPLEMENTED and automated (tag-triggered). \\
Extension delivery & NOT AUTOMATED; artefact committed by hand. \\
Environments & Single (production); no staging or preview. \\
Containerisation / IaC & NOT APPLICABLE. \\
Rollback & INFERRED only; not documented. \\
\bottomrule
\end{tabular}
\end{table}

\textbf{Status:} IMPLEMENTED (registry, npm) / NOT AUTOMATED (extension) /
RECOMMENDED (extension release automation, rollback documentation).
